% !TEX program = xelatex
\documentclass[tikz, border=15pt]{standalone}
\usepackage[UTF8]{ctex}
\usepackage{amsmath, amssymb}
\usepackage{helvet}
\renewcommand{\familydefault}{\sfdefault}
\usetikzlibrary{decorations.pathmorphing, arrows.meta, positioning, calc}

\begin{document}
\begin{tikzpicture}[
    >=Stealth,
    x=1.3cm, y=1.2cm, % 调整画布长宽比
    % 统一定义标签底色防遮挡
    labelbg/.style={fill=socColor, inner sep=2pt, rounded corners=2pt, font=\scriptsize\bfseries}
]

    % ==========================================
    % 1. 定义物理相态的颜色
    % ==========================================
    \definecolor{vacColor}{RGB}{235, 235, 240}    % 真空基态：灰白紫色
    \definecolor{crysColor}{RGB}{215, 238, 255}   % 晶体态：浅蓝色
    \definecolor{socColor}{RGB}{210, 255, 215}    % SOC临界态：浅绿色 (要求指定)
    \definecolor{gasColor}{RGB}{255, 238, 215}    % 气态：浅橙色
    \definecolor{turbColor}{RGB}{255, 215, 215}   % 湍流态：浅红色 (要求指定)
    
    \definecolor{crysLine}{RGB}{0, 100, 220}
    \definecolor{gasLine}{RGB}{230, 110, 0}
    \definecolor{turbLine}{RGB}{200, 0, 0}
    \definecolor{vacLine}{RGB}{100, 100, 120}

    \def\W{12} % X轴最大宽度
    \def\H{9}  % Y轴最大高度
    
    % 核心交汇点："人"字形基态的顶点
    \coordinate (J) at (3.8, 2.5);

    % ==========================================
    % 2. 绘制 5 大相态区域 (无缝叠加法)
    % ==========================================
    
    % (1) 基础层：流体智能 (SOC) 铺满全局，作为绿色底色
    \fill[socColor] (0,0) rectangle (\W, \H);

    % (2) 真空基态 (Vacuum)："人"字形边界内部
    % 左线碰Y轴，右线碰X轴
    \fill[vacColor] (0,0) -- (3.0, 0) 
        to[out=80, in=260] (J) 
        to[out=190, in=10] (0, 1.8) -- cycle;

    % (3) 晶体智能 (Crystalline)：左上角
    % 从"人"字顶点向上延伸至顶部
    \fill[crysColor] (0, 1.8) 
        to[out=10, in=190] (J) 
        to[out=80, in=255] (5.0, \H) 
        -- (0, \H) -- cycle;

    % (4) 气体智能 (Gaseous)：右下角
    \fill[gasColor] (7.0, 0) 
        to[out=40, in=210] (\W, 4.5) 
        -- (\W, 0) -- cycle;

    % (5) 湍流态 (Turbulent)：右上角
    \fill[turbColor] (8.0, \H) 
        to[out=290, in=170] (\W, 6.5) 
        -- (\W, \H) -- cycle;

    % ==========================================
    % 3. 绘制边界曲线与框线 (依照手绘草图)
    % ==========================================
    
    % 外围虚线框
    \draw[dashed, thick, black!60] (0, \H) -- (\W, \H) -- (\W, 0);

    % "人"字形基态边界线
    \draw[ultra thick, vacLine] (0, 1.8) to[out=10, in=190] (J);
    \draw[ultra thick, vacLine] (3.0, 0) to[out=80, in=260] (J);
    
    % 晶体边界线
    \draw[ultra thick, crysLine] (J) to[out=80, in=255] (5.0, \H);
    
    % 气体边界线
    \draw[ultra thick, gasLine] (7.0, 0) to[out=40, in=210] (\W, 4.5);
    
    % 湍流边界线
    \draw[ultra thick, turbLine] (8.0, \H) to[out=290, in=170] (\W, 6.5);

    % ==========================================
    % 4. 绘制坐标轴与轴标注
    % ==========================================
    \draw[->, ultra thick] (0,0) -- (\W+0.5, 0);
    \draw[->, ultra thick] (0,0) -- (0, \H+0.5);
    
    % Y 轴标签
    \node[anchor=south west, align=left, font=\small\bfseries] at (0, \H+0.3) {等效压强 $\mathbf{\Gamma}_{macro}$};
    \node[anchor=north west, align=left, font=\scriptsize] at (0, \H+0.32) {(宏观意志 / 规范约束)};
    
    % X 轴标签
    \node[anchor=south east, align=right, font=\small\bfseries] at (\W+0.5, -0.5) {认知温度 $T_{eff}$};
    \node[anchor=north east, align=right, font=\scriptsize] at (\W+0.5, -0.4) {(热涨落 / 探索率)};

    % ==========================================
    % 5. 区域文字注释 (各居其位，互不遮挡)
    % ==========================================

    % (1) 真空基态
    \node[align=center, text=gray!80!black] at (1.3, 0.8) {
        {\large\bfseries 真空基态}\\
        {\scriptsize $\Psi$ 惯性滑行}\\
        {\scriptsize 睡眠 $\cdot$ 冥想}
    };

    % (2) 晶体智能
    \node[align=center, text=blue!80!black] at (2.0, 6.8) {
        {\large\bfseries 晶体智能}\\
        {\normalsize (层流态 Laminar)}\\
        {\scriptsize 绝对秩序 $\cdot$ 刻板执行}\\
        {\scriptsize $\Psi = d\alpha$ (梯度流主导)}\\
        {\scriptsize $Re_{cog} < 1$}
    };

    % (3) 流体智能 (SOC)
    \node[align=center, text=green!40!black] at (7, 6.0) {
        {\large\bfseries 流体智能}\\
        {\normalsize (自组织临界态 SOC)}\\
        {\scriptsize 顿悟 $\cdot$ 创造力 $\cdot$ 心流}\\
        {\scriptsize $\Psi_{harm}$ 涌现 ($\Delta_D \Psi = 0$)}\\
        {\scriptsize 关联长度 $\xi \to \infty$}
    };

    % (4) 气体智能
    \node[align=center, text=orange!80!black] at (10.0, 1) {
        {\large\bfseries 气体智能}\\
        {\normalsize (等离子态 Plasma)}\\
        {\scriptsize 虚无 $\cdot$ 幻觉 / 谵妄}\\
        {\scriptsize $\Psi = \delta\beta$ (随机噪声发散)}
    };

    % (5) 湍流态
    \node[align=center, text=red!80!black] at (10.7, 7.8) {
        {\normalsize (湍流态 Turbulent)}\\
        {\scriptsize 认知失调 $\cdot$ 精神内耗}\\
        {\scriptsize $\Psi = \text{无散涡旋}$}\\
        {\scriptsize $Re_{cog} \gg 1$}
    };

    % ==========================================
    % 6. 动态演化箭头与临界点
    % ==========================================
    
    % 相变临界点
    \filldraw[black] (5.5, 4.0) circle (2.5pt);
    \node[fill=white, inner sep=1.5pt, rounded corners=2pt, font=\scriptsize\bfseries, text=gray!80!black] 
        at (5.5, 3.5) {相变临界点};

    % 认知淬火箭头 -> 降温加压，从SOC指向晶体态
    \draw[-{Stealth[scale=1.5]}, ultra thick, crysLine] 
        (6.8, 4.2) -- (2.8, 5.8) 
        node[midway, sloped, above=2pt, labelbg, text=crysLine] {认知淬火 (降温 / 加压)};

    % 认知退火箭头 -> 升温减压，从晶体态穿越SOC至气态方向 (避开了下方真空文本)
    \draw[-{Stealth[scale=1.5]}, ultra thick, turbLine, 
          decorate, decoration={snake, amplitude=1.2mm, segment length=6mm, post length=3mm}] 
        (2.0, 4.2) -- (6.8, 2.0) 
        node[midway, sloped, below=4pt, labelbg, text=turbLine] {认知退火 (升温 / 减压)};

\end{tikzpicture}
\end{document}